\section{Discussione e Conclusioni}

L'ablazione mostra che la riduzione dei cicli dipende dalla ricerca TOP-4.
Il classificatore può scartare dati utili.
TOP-4 è efficace nel caso studiato, ma non è una scelta ottima per ogni istanza.

Nella correzione d'errore, la rete aiuta il matching quando sono presenti
interferenze fra qubit. Il beneficio diminuisce con la distanza del codice.

Tutti gli esperimenti sono simulazioni. Lo studio del classificatore usa
solo $N=15$, con rumore uniforme e invariato nel tempo.
Per il classificatore, gli errori di operazioni diverse sono indipendenti; non si considerano
interferenze fra qubit né operazioni aggiuntive richieste dalle loro connessioni.

La correzione considera soltanto alcune operazioni e alcuni tipi di errore:
le stime possono quindi essere favorevoli~\cite{plaquette2026}.
Non viene simulato un circuito completo di Shor protetto dai codici.

Il vantaggio della QFT approssimata nella stima di fase non si estende
automaticamente a Shor. Le prove rumorose su $N=21$ confrontano varianti
della QFT; $N=35$ è verificato senza rumore.
La stima di fase usa una disposizione dei qubit per dimensione,
tre fasi e un solo insieme di dati di calibrazione.

Occorre collegare gli errori dopo la correzione alle operazioni di Shor
ed estendere i modelli di rumore.
Per i decoder, si potranno usare reti che considerino disposizione dei
qubit ed evoluzione dei controlli~\cite{alphaqubit2024}.
La QFT approssimata andrà combinata con metodi aritmetici che riducono
i qubit necessari~\cite{chevignard2025}, confrontando porte, qubit e profondità.
